\section{Complementary Filter for Attitude Estimation}
\label{sec:complementary_filter}

% ----------------------------------------------------------------
% TARGET: ~1.5 pages. The core fusion algorithm.
% ----------------------------------------------------------------

\subsection{Intuition}
\label{sec:cf:intuition}

As discussed in Section~\ref{sec:frames:fusion_need}, the accelerometer and
gyroscope have complementary noise characteristics. The key insight is that
their useful information occupies \emph{different frequency bands}:

\begin{itemize}
  \item The accelerometer-derived angle is reliable at \textbf{low frequencies}
        (below $\sim$\!1~Hz) but corrupted at high frequencies.
  \item The gyroscope-integrated angle is reliable at \textbf{high frequencies}
        but drifts at low frequencies.
\end{itemize}

A complementary filter combines the two by applying a \emph{low-pass filter}
to the accelerometer and a \emph{high-pass filter} to the gyroscope, such
that their transfer functions sum to unity across all frequencies.

\subsection{Filter Derivation}
\label{sec:cf:derivation}

Let $\theta_a(t)$ denote the tilt angle estimated from the accelerometer and
$\theta_g(t)$ the angle obtained by integrating the gyroscope. The
complementary filter output is:

\begin{equation}
  \boxed{
    \hat{\theta}(t) = \alpha \bigl[\hat{\theta}(t - \Delta t)
    + \omega(t)\,\Delta t\bigr]
    + (1 - \alpha)\,\theta_a(t)
  }
  \label{eq:comp_filter}
\end{equation}

\noindent where $\alpha \in (0, 1)$ is the filter coefficient. The first term
is a \emph{high-pass filter} on the gyroscope (it propagates the previous
estimate forward and adds the incremental rotation), and the second term is a
\emph{low-pass filter} on the accelerometer (it slowly corrects toward the
gravity reference).

The time constant of the filter is:

\begin{equation}
  \tau = \frac{\alpha\,\Delta t}{1 - \alpha}
  \label{eq:time_constant}
\end{equation}

\noindent Typical values are $\alpha = 0.96$--$0.98$, corresponding to time
constants of a few seconds.

% ----- FIGURE: complementary filter block diagram -----
% TODO: Block diagram showing:
%   - Gyroscope → integrator → HPF ──┐
%   - Accelerometer → atan2 → LPF ──┤→ Σ → θ_hat
%   - Labels: α and (1-α)
%   - Frequency response subplot below
%
% \begin{figure}[t]
%   \centering
%   \includegraphics[width=\columnwidth]{complementary_filter.pdf}
%   \caption{Block diagram of the complementary filter. The gyroscope path
%            (top) acts as a high-pass filter with coefficient $\alpha$; the
%            accelerometer path (bottom) acts as a low-pass filter with
%            coefficient $(1-\alpha)$. Their sum reconstructs the full-band
%            attitude estimate.}
%   \label{fig:comp_filter}
% \end{figure}

\subsection{Accelerometer Tilt Extraction}
\label{sec:cf:tilt}

When the sensor is approximately stationary, the accelerometer measures the
gravity vector projected onto the body axes. The roll and pitch angles can be
extracted as:

\begin{align}
  \phi   &= \arctan\!\left(\frac{a_y}{\sqrt{a_x^2 + a_z^2}}\right)
  \label{eq:roll} \\[4pt]
  \theta &= \arctan\!\left(\frac{-a_x}{\sqrt{a_y^2 + a_z^2}}\right)
  \label{eq:pitch}
\end{align}

\noindent These relations assume that the only force acting on the sensor is
gravity. During dynamic motion, the estimates become unreliable---precisely the
regime where the gyroscope takes over through the high-pass path.

\subsection{Quaternion Form of the Complementary Filter}
\label{sec:cf:quaternion}

To avoid gimbal lock, the complementary filter is implemented in quaternion
space. At each time step:

\begin{enumerate}
  \item Integrate the gyroscope to obtain a predicted quaternion
        $\mathbf{q}_\text{pred}$ using Eq.~\eqref{eq:quat_euler}.
  \item Compute a reference quaternion $\mathbf{q}_\text{acc}$ from the
        accelerometer tilt angles (Eqs.~\ref{eq:roll}--\ref{eq:pitch}).
  \item Blend the two using spherical linear interpolation (SLERP):
\end{enumerate}

\begin{equation}
  \mathbf{q}(t) = \text{SLERP}\bigl(\mathbf{q}_\text{acc},\;
                  \mathbf{q}_\text{pred},\; \alpha\bigr)
  \label{eq:slerp_filter}
\end{equation}

\noindent This is the quaternion analogue of the scalar
Eq.~\eqref{eq:comp_filter}, with the SLERP parameter playing the role of
$\alpha$.
